                  6th INTERNATIONAL COMPETITION FOR UNIVERSITY
                             STUDENTS IN MATHEMATICS
                                    Keszthely, 1999.

                                       Problems and solutions on the first day

1. a) Show that for any m ∈ N there exists a real m × m matrix A such that A3 = A + I, where I is the
m × m identity matrix. (6 points)
b) Show that det A > 0 for every real m × m matrix satisfying A3 = A + I. (14 points)
Solution. a) The diagonal matrix                                                  
                                                                 λ             0
                                                                      ..
                                               A = λI =                   .       
                                                                 0             λ
is a solution for equation A3 = A + I if and only if λ3 = λ + 1, because A3 − A − I = (λ3 − λ − 1)I. This
equation, being cubic, has real solution.
      b) It is easy to check that the polynomial p(x) = x3 − x − 1 has a positive real root λ1 (because p(0) < 0)
and two conjugated complex roots λ2 and λ3 (one can check the discriminant of the polynomial, which is
     
  −1 3
                2    23
   3    + −1  2    = 108 > 0, or the local minimum and maximum of the polynomial).
      If a matrix A satisfies equation A3 = A + I, then its eigenvalues can be only λ1 , λ2 and λ3 . The
multiplicity of λ2 and λ3 must be the same, because A is a real matrix and its characteristic polynomial has
only real coefficients. Denoting the multiplicity of λ1 by α and the common multiplicity of λ2 and λ3 by β,
                                                        β β
                                            det A = λα          α           β
                                                     1 λ2 λ3 = λ1 · (λ2 λ3 ) .

Because λ1 and λ2 λ3 = |λ2 |2 are positive, the product on the right side has only positive factors.

2. Does there exist a bijective map π : N → N such that
                                                     ∞
                                                     X π(n)
                                                                     < ∞?
                                                     n=1
                                                           n2

(20 points)
Solution 1. No. For, let π be a permutation of N and let N ∈ N. We shall argue that

                                                      3N
                                                      X  π(n)  1
                                                            2
                                                              > .
                                                          n    9
                                                    n=N +1


In fact, of the 2N numbers π(N + 1), . . . , π(3N ) only N can be ≤ N so that at least N of them are > N .
Hence
                           3N                        3N
                           X   π(n)           1     X             1           1
                                  2
                                      ≥           2
                                                         π(n) >     2
                                                                      ·N ·N = .
                                n           (3N )               9N            9
                             n=N +1                     n=N +1


Solution 2. Let π be a permutation of N. For any n ∈ N, the numbers π(1), . . . , π(n) are distinct positive
integers, thus π(1) + . . . + π(n) ≥ 1 + . . . + n = n(n+1)
                                                        2   . By this inequality,
                            ∞               ∞                                                     
                            X π(n)          X                                     1       1
                                        =         π(1) + . . . + π(n)                 −                ≥
                            n=1
                                  n2        n=1
                                                                                   n2   (n + 1)2

                           ∞                               ∞               ∞
                           X n(n + 1)          2n + 1     X    2n + 1     X     1
                       ≥                    · 2         =               ≥           = ∞.
                           n=1
                                  2          n (n + 1)2   n=1
                                                              2n(n + 1)   n=1
                                                                              n + 1

                                                             1
3. Suppose that a function f : R → R satisfies the inequality
                                         n
                                         X                                   
                                               3k f (x + ky) − f (x − ky)        ≤1                               (1)
                                         k=1

for every positive integer n and for all x, y ∈ R. Prove that f is a constant function. (20 points)
Solution. Writing (1) with n − 1 instead of n,
                                        n−1
                                        X                                    
                                               3k f (x + ky) − f (x − ky)        ≤ 1.                             (2)
                                        k=1

From the difference of (1) and (2),
                                                                         
                                          3n f (x + ny) − f (x − ny)         ≤ 2;

which means
                                                                2
                                                                  .
                                            f (x + ny) − f (x − ny) ≤                               (3)
                                                               3n
    For arbitrary u, v ∈ R and n ∈ N one can choose x and y such that x − ny = u and x + ny = v, namely
x = u+v           v−u
     2 and y = 2n . Thus, (3) yields
                                                           2
                                           f (u) − f (v) ≤ n
                                                          3
for arbitrary positive integer n. Because 32n can be arbitrary small, this implies f (u) = f (v).
                                                                              2 
4. Find all strictly monotonic functions f : (0, +∞) → (0, +∞) such that f fx(x) ≡ x. (20 points)
                          f (x)               x                                              x
Solution. Let g(x) =            . We have g(      ) = g(x). By induction it follows that g( n    ) = g(x), i.e.
                            x                g(x)                                          g (x)
                                                x         x
(1)                                         f( n    ) = n−1     , n ∈ N.
                                              g (x)    g    (x)

                                                    x2                                             f 2 (x)
      On the other hand, let substitute x by f (x) in f () = x. ¿From the injectivity of f we get           =
                                                   f (x)                                          f (f (x))
x, i.e. g(xg(x)) = g(x). Again by induction we deduce that g(xg n (x)) = g(x) which can be written in the
form

(2)                                        f (xg n (x)) = xg n−1 (x), n ∈ N.

      Set f (m) = f ◦ f ◦ . . . ◦ f . It follows from (1) and (2) that
                  |     {z        }
                      m times


(3)                                    f (m) (xg n (x)) = xg n−m (x), m, n ∈ N.

     Now, we shall prove that g is a constant. Assume g(x1 ) < g(x2 ). Then we may find n ∈ N such
that x1 g n (x1 ) ≤ x2 g n (x2 ). On the other hand, if m is even then f (m) is strictly increasing and from (3) it
follows that xm  1 g
                     n−m
                         (x1 ) ≤ xm2 g
                                       n−m
                                           (x2 ). But when n is fixed the opposite inequality holds ∀m  1. This
contradiction shows that g is a constant, i.e. f (x) = Cx, C > 0.
     Conversely, it is easy to check that the functions of this type verify the conditions of the problem.

5. Suppose that 2n points of an n × n grid are marked. Show that for some k > 1 one can select 2k distinct
marked points, say a1 , . . . , a2k , such that a1 and a2 are in the same row, a2 and a3 are in the same column,
. . . , a2k−1 and a2k are in the same row, and a2k and a1 are in the same column. (20 points)

                                                            2
Solution 1. We prove the more general statement that if at least n + k points are marked in an n × k grid,
then the required sequence of marked points can be selected.
     If a row or a column contains at most one marked point, delete it. This decreases n + k by 1 and the
number of the marked points by at most 1, so the condition remains true. Repeat this step until each row
and column contains at least two marked points. Note that the condition implies that there are at least two
marked points, so the whole set of marked points cannot be deleted.
     We define a sequence b1 , b2 , . . . of marked points. Let b1 be an arbitrary marked point. For any positive
integer n, let b2n be an other marked point in the row of b2n−1 and b2n+1 be an other marked point in the
column of b2n .




      Let m be the first index for which bm is the same as one of the earlier points, say bm = bl , l < m.
      If m − l is even, the line segments bl bl+1 , bl+1 bl+2 , ..., bm−1 bl = bm−1 bm are alternating horizontal and
vertical. So one can choose 2k = m − l, and (a1 , . . . , a2k ) = (bl , . . . , bm−1 ) or (a1 , . . . , a2k ) = (bl+1 , . . . , bm )
if l is odd or even, respectively.
      If m − l is odd, then the points bl = bm , bl+1 and bm−1 are in the same row/column. In this case chose
2k = m − l − 1. Again, the line segments bl+1 bl+2 , bl+2 bl+3 , ..., bm−1 bl+1 are alternating horizontal and
vertical and one can choose (a1 , . . . , a2k ) = (bl+1 , . . . , bm−1 ) or (a1 , . . . , a2k ) = (bl+2 , . . . , bm−1 , bl+1 ) if l is
even or odd, respectively.
Solution 2. Define the graph G in the following way: Let the vertices of G be the rows and the columns of
the grid. Connect a row r and a column c with an edge if the intersection point of r and c is marked.
     The graph G has 2n vertices and 2n edges. As is well known, if a graph of N vertices contains no circle,
it can have at most N − 1 edges. Thus G does contain a circle. A circle is an alternating sequence of rows
and columns, and the intersection of each neighbouring row and column is a marked point. The required
sequence consists of these intersection points.

6. a) For each 1 < p < ∞ find a constant cp < ∞ for which the following statement holds: If f : [−1, 1] → R
is a continuously differentiable function satisfying f (1) > f (−1) and |f 0 (y)| ≤ 1 for all y ∈ [−1, 1], then there
                                                                          1/p
is an x ∈ [−1, 1] such that f 0 (x) > 0 and |f (y) − f (x)| ≤ cp f 0 (x)       |y − x| for all y ∈ [−1, 1]. (10 points)
b) Does such a constant also exist for p = 1? (10 points)
                                                             R1                 R1             R1
Solution. (a) Let g(x) = max(0, f 0 (x)). Then 0 < −1 f 0 (x)dx = −1 g(x)dx + −1 (f 0 (x) − g(x))dx, so
         R1                 R1             R1                          R1
we get −1 |f 0 (x)|dx = −1 g(x)dx + −1 (g(x) − f 0 (x))dx < 2 −1 g(x)dx. Fix p and c (to be determined
at the end). Given any t > 0, choose for every x such that g(x) > t an interval I x = [x, y] such that
|f (y) − f (x)| > cg(x)1/p |y − x| > ct1/p |Ix | and choose disjoint Ixi that cover at least one third of the measure
                                   S                                 R               R1                  R1
of the set {g > t}. For I = i Ii we thus have ct1/p |I| ≤ I f 0 (x)dx ≤ −1 |f 0 (x)|dx < 2 −1 g(x)dx; so
                                  R 1                                                  R 1            R 1
|{g > t}| ≤ 3|I| < (6/c)t−1/p −1 g(x)dx. Integrating the inequality, we get −1 g(x)dx = 0 |{g > t}|dt <
                 R1
(6/c)p/(p − 1) −1 g(x)dx; this is a contradiction e.g. for cp = (6p)/(p − 1).
      (b) No. Given c > 1, denote α = 1/c and choose 0 < ε < 1 such that ((1 + ε)/(2ε)) −α < 1/4. Let
g : [−1, 1] → [−1, 1] be continuous, even, g(x) = −1 for |x| ≤ ε and 0 ≤ g(x) < α((|x| + ε)/(2ε)) −α−1 for ε <
                              R1                    R1
|x| ≤ 1 is chosen such that ε g(t)dt > −ε/2+ ε α((|x|+ε)/(2ε))−α−1 dt = −ε/2+2ε(1−((1+ε)/(2ε))−α) > ε.
          R                                            R1
Let f = g(t)dt. Then f (1) − f (−1) ≥ −2ε + 2 ε g(t)dt > 0. If ε < x < 1 and y = −ε, then |f (x) − f (y)| ≥
       Rx                  Rx
2ε − ε g(t)dt ≥ 2ε − ε α((t + ε)/(2ε))−α−1 = 2ε((x + ε)/(2ε))−α > g(x)|x − y|/α = f 0 (x)|x − y|/α;
symmetrically for −1 < x < −ε and y = ε.



                                                                  3
